\documentclass[11pt,a4paper]{article}
\usepackage{fontspec,amsmath,amssymb,amsthm,microtype}
\setmainfont{lmroman10-regular.otf}[BoldFont=lmroman10-bold.otf,ItalicFont=lmroman10-italic.otf,BoldItalicFont=lmroman10-bolditalic.otf]
\setmonofont{lmmono10-regular.otf}[Scale=MatchLowercase]
\usepackage[margin=25mm]{geometry}
\usepackage[unicode,hidelinks]{hyperref}
\setlength{\emergencystretch}{3em}
\setlength{\parskip}{0.4em}
\setlength{\parindent}{0pt}
\allowdisplaybreaks[2]
\raggedbottom
\title{\texorpdfstring{Compatibility exports for constructed Brownian reference segments}{Compatibility exports for constructed Brownian reference segments}}
\author{LangevinLean mathematical exposition}
\date{}
\begin{document}\maketitle
\noindent Source module: \texttt{IPLA/\allowbreak{}Brownian\allowbreak{}Reference.\allowbreak{}lean}. The accompanying HTML pairs this mathematical exposition with the exact Lean statements and records the scope of the formalization.


This file exports shared Brownian segment and reference constructions and contains no duplicated local declarations. The imported Brownian-reference layer has an explicit all-time Gibbs-law premise; the complete Gaussian ULA development later proves stationarity.


\section{\texorpdfstring{Join integral solutions and identify their Euler noises}{Join integral solutions and identify their Euler noises}}


For each interval \([kh,(k+1)h]\), solve the Lipschitz-drift integral equation in \(L^2\), starting from the preceding endpoint:

\[Y_{k,t}-Y_{k,0}=-\int_0^t\nabla V(Y_{k,s})ds+\sqrt2\bigl(B_{kh+t}-B_{kh}\bigr).\]

The construction proves continuity, the initial value and interval joining. Normalize each Brownian endpoint by \(\sqrt{2h}\) to obtain the same Euler noise vector, and use Gaussian second moments to get the noise norm \(\sqrt{2d}\sqrt t\). Given the explicit all-time Gibbs-law premise, the Gibbs score bound yields \(\|\nabla V(Y_{k,t})\|_{L^2}\le\sqrt{Ld}\), assembling the shared integral-reference structure. This layer keeps stationarity visible as a premise, rather than claiming the construction alone proves it.



The \texttt{export} commands expose the already proved declarations under the IPLA namespace. Their definitions and proofs are unchanged, so the mathematical argument is the one in the imported source modules; the dependency diagram below opens those full proofs. This file itself makes no additional analytic assumption and contains no replacement proof.


\end{document}
